\section{Introduction}
\label{sec:introduction}

A robot can encounter the same task description at several stages of an interaction, each requiring a different next decision. To place a cup in a drawer, it may first need to open the drawer, recover from a failed grasp, or verify that the cup was released inside. Relevant experience depends on the current phase and on what previous actions actually achieved. A record of successful placement is useful only if its prerequisites can be satisfied in the present scene. A failed attempt can be equally useful when its feedback explains why a different observation or action was needed.

Consider a robot with experience transferring objects and opening drawers, but no recorded episode of placing a cup into a closed drawer. Solving this new composition requires more than matching the instruction to a familiar task. For a single-gripper robot, drawer opening and handle release must precede the cup grasp. The transfer procedure must bind its object and destination roles to the current cup and drawer, and ground motion parameters from current observations. If the grasp fails, the robot must identify which condition remains unmet and follow an appropriate recovery. This example connects three questions: how should recorded interactions become reusable experience, how should that experience change a new plan, and how should the robot determine whether the plan is being fulfilled?

Retrieval and structured planning offer substantial foundations for answering these questions. Memory Transfer Planning retrieves and adapts successful planner code for new manipulation tasks~\cite{mtp2025}. HP-KG constructs hierarchical procedural knowledge and retrieves relevant subgraphs for planning~\cite{hpkg2025}. Neural Task Graphs connects demonstration-derived task graphs to observation-conditioned execution~\cite{ntg2019}, while GAVEL uses a graph world model to validate and repair plans~\cite{gavel2026}. These results motivate a more specific research question: can preserving the relations among phase conditions, attempted actions, observed outcomes, and subsequent choices improve how experience is selected and used throughout manipulation?

We introduce \emph{RoboRAG: Interaction-Aware Retrieval-Augmented Manipulation}, a framework organized around this question (Fig.~\ref{fig:overview}). Its central hypothesis is that explicit conditions and feedback links make retrieved experience more applicable to the current decision and make plan progress easier to verify. A shared phase-graph schema connects historical experience, plans, and execution, while current evidence determines whether planned effects occurred.

The first method component constructs the graph bank through an annotation pipeline. It segments timestamped interactions into phases, assigns object roles and local objectives, and links actions to feedback and evidence. Nodes represent phase instances; typed edges express precedence, prerequisites, conditional transitions, and recovery. The resulting graph preserves the timing of attempts and the reason a continuation became relevant. Structured records support validation, while compact descriptions expose the same semantics to a language model. Whole-graph and phase-centered indexes allow retrieval of either a complete procedure or a smaller pattern with the conditions needed to interpret it. Graph versions retain the provenance of corrections and additions.

The second component retrieves, adapts, and follows these graphs. A GPT planner constructs a provisional graph for an unseen task using a declared schema and backend capabilities. Retrieval finds compatible procedures and local experience; adaptation aligns roles, grounds parameters, and edits dependencies or recovery branches. In the drawer example, opening experience supplies the missing prerequisite for the transfer plan. A lower-cost supervisor then selects eligible phases and approved recovery steps. A deterministic gateway checks requests, and native controllers execute skills. Structural revisions return to the planner. This hierarchy makes the plan an operational contract with observable completion conditions and gives routine execution a defined path without repeated high-level planning.

The manuscript develops three contributions:
\begin{enumerate}
\item An interaction-graph construction and indexing method that retains phase boundaries, conditions, observed feedback, and recovery relations in reusable experience.
\item A graph retrieval, adaptation, and execution method that carries those relations into plans for unseen task compositions and uses current evidence to govern delegated plan following.
\item An evaluation design that isolates the value of these relations through matched evidence and execution comparisons, together with a common application contract for studying portability across RoboDojo, LIBERO, and a real robot.
\end{enumerate}

The evaluation asks whether graph structure improves decisions beyond linearized versions of the same evidence, whether adaptation supports unseen compositions, and whether delegated execution preserves task performance while reducing high-level calls. A separate systems comparison distinguishes backend normalization from MCP transport. Together, these tests connect the proposed representation to observable planning, recovery, execution, and integration outcomes.

% CHAPTER_SUMMARY: RoboRAG asks whether retaining phase conditions, observed outcomes, and recovery relations across experience construction, retrieval/adaptation, and delegated execution improves manipulation of unseen task compositions. A cup-in-closed-drawer example motivates a shared schema whose historical, planned, and executed graphs retain distinct evidence meanings. Two method components build and use the graph bank; matched experiments test the representation, while the cross-backend common contract and MCP comparison support systems validation.
